\documentclass[11pt]{article}

\usepackage[a4paper,margin=1in]{geometry}
\usepackage{amsmath,amssymb,amsthm,mathtools}
\usepackage[hidelinks]{hyperref}
\hypersetup{
  pdfauthor={Bangyen Pham},
  pdftitle={The Worst-Case Text Complexity of the Polynomial Esoteric Programming Language},
  pdfsubject={Descriptional complexity of Boolean functions in Polynomial},
  pdfkeywords={Polynomial language, descriptional complexity, Boolean functions, decision diagrams, polynomial roots},
  pdflang={en-US}
}

\newtheorem{theorem}{Theorem}[section]
\newtheorem{lemma}[theorem]{Lemma}
\newtheorem{corollary}[theorem]{Corollary}
\newtheorem{proposition}[theorem]{Proposition}
\theoremstyle{remark}
\newtheorem{remark}[theorem]{Remark}

\newcommand{\R}{\mathbb R}
\newcommand{\Z}{\mathbb Z}
\newcommand{\mroute}{m_{\mathrm{route}}}

\title{The Worst-Case Text Complexity of the Polynomial Language}
\author{\shortstack{Bangyen Pham\\
\small Independent Researcher\\
\small\href{mailto:bangyenp@gmail.com}{\texttt{bangyenp@gmail.com}}\\
\small\href{https://orcid.org/0009-0006-9928-2088}{ORCID: 0009-0006-9928-2088}}}
\date{September 2026}

\begin{document}
\maketitle

\begin{abstract}
In the Polynomial esoteric programming language, the roots of an
integer polynomial encode the instructions of a one-register machine, while
the source text lists the polynomial's expanded coefficients.  We prove that
every Boolean function with truth-table length $T$ has a Polynomial program of
length $O(T^2/\log T)$, and that some require
$\Omega(T^2/\log T)$ source characters in every program, independently of the program's degree, operand signs, and additional
roots.
The explicit residual-DAG construction gives the upper bound.  A
coefficient-mass theorem, proved separately, implies that every real
polynomial divisible by $L$ distinct prime-power roots and having leading
coefficient of absolute value at least one has logarithmic coefficient mass
$\Omega(L^2\log L)$.  A routing argument, sharpened by an
outerplanar block-incidence bound, forces $L=\Omega(T/\log T)$ distinct real
roots in a program for a maximal-width table, so the worst-case text
complexity is $\Theta(T^2/\log T)$ in this model.  The routing argument
needs no assumption on the number of reads: a read cursor that routing
cannot distinguish carries a residual that ignores its next input.
\end{abstract}

\noindent\textbf{Keywords.} Descriptional complexity; Boolean functions;
ordered binary decision diagrams; polynomial roots; esoteric programming
languages.

\noindent\textbf{2020 Mathematics Subject Classification.} Primary 68Q17;
secondary 94D10, 11C08.

\section{The language and the result}\label{sec:model}

A Polynomial source is the expanded form
\[
  f(x)=\sum_{j=0}^{D} f_jx^j,\qquad f_j\in\Z.
\]
It is serialized as \texttt{f(x) = } followed by nonzero monomials in
descending degree, with decimal coefficients and exponents, as in the
language specification~\cite{polynomial-spec}; whitespace is permitted
between tokens.  Throughout, ``program'' means a source in this specified
expanded form; permissive implementation extensions are outside the model.
The interpreter factors $f$ over the integers and orders the recognized
roots by their associated primes.  A positive real root $p^v$ encodes a
control-flow instruction; a conjugate pair $a\pm p^b i$ encodes a register
instruction.  Unrecognized roots are ignored, so a source may be any integer
multiple of the product of its instruction factors.  Repeated recognized
roots produce repeated instructions.

As a minimal example, $f(x)=x^2+4$ has roots $\pm2i$.  The imaginary
magnitude $2=2^1$ selects the first register instruction for the prime $2$,
while the real part $0$ supplies its operand.  This instruction outputs the
initial register value, the zero byte.  In general the interpreter follows
the same pipeline:
\[
 \text{expanded coefficients}\longrightarrow
 \text{recognized roots}\longrightarrow
 \text{sorted instructions}\longrightarrow
 \text{execution}.
\]

The machine has an integer register and an instruction cursor.  Its input
instruction overwrites the register with the next input symbol, $48+b$ for a
Boolean bit $b$.  Real-root instructions are the bracket-like control-flow
operations; complex-root instructions update the register.  The upper bound
uses input, output, addition and multiplication by an integer, integer floor
division, and ``if positive'' with its closing instruction.  The lower bound
uses only the following structural properties.

\begin{enumerate}
\item Input is consumed as an ordered stream: the $j$-th read instruction
receives the $j$-th input bit, whatever the path taken.
\item Between two input reads the complete state is the pair consisting of
the register and the instruction cursor.
\item Each real instruction root is one of the first eight positive powers of
its associated prime.  One prime may supply several consecutive positions.
\item Equal real roots occur in one contiguous block after root sorting.
\item The bracket matching is noncrossing.
\end{enumerate}

For Boolean computation, put $T=2^n$ and supply each bit as a separate line
containing its ASCII digit.  A program computes
$t=(t_0,\ldots,t_{T-1})\in\{0,1\}^{2^n}$ if, on every
$b_1\cdots b_n\in\{0,1\}^n$ in lexicographic order, it terminates and emits
exactly the character $t_{(b_1\cdots b_n)_2}$, with zero-based subscripts.
Reads consume the input stream in order; a read past its end stores $-1$.
A program may perform different numbers of reads on different inputs, and
may read past the end.  The upper-bound construction reads exactly $n$ bits
on every input, but the lower bound does not assume this.

The length $|f|$ is the total
number of characters in the shortest valid expanded serialization of $f$;
fixed punctuation and optional whitespace affect only lower-order terms.
Its coefficient-digit mass is
\[
  M(f)=\sum_{j:f_j\ne0}\bigl(1+\lfloor\log_{10}|f_j|\rfloor\bigr).
\]
Every nonzero coefficient appears in a monomial whose serialization supplies
at least $1+\lfloor\log_{10}|f_j|\rfloor$ characters, so $|f|\ge M(f)$.
The worst-case text complexity is the maximum, over these tables, of the
minimum $|f|$ of a program computing the table.
All unqualified logarithms are natural.

\begin{theorem}[Main theorem]\label{thm:main}
Every $n$-input Boolean function has a Polynomial program of length
$O(T^2/\log T)$, and there are $n$-input Boolean functions for which every
program has length $\Omega(T^2/\log T)$.  Hence the
worst-case text complexity in the model of Section~\ref{sec:model} is
$\Theta(T^2/\log T)$.
\end{theorem}

The lower bound has two independent parts.
Section~\ref{sec:routing} forces $\Omega(T/\log T)$ distinct positive
roots.
Section~\ref{sec:mass} applies the
coefficient-mass theorem of~\cite{coefficient-mass}.  Section~\ref{sec:upper} supplies
the matching construction.  Informally, the routing argument says that a
short instruction alphabet cannot distinguish enough residual truth tables;
the coefficient-mass theorem then says that supporting the required roots
forces many coefficient digits.

Residual DAGs are reduced ordered binary decision diagrams under a fixed
variable order~\cite{bryant}; classical results place the size of an optimal
diagram for almost every Boolean function at the same
$\Theta(T/\log T)$ scale~\cite{wegener,gps}.  Lemma~\ref{lem:maximal-width}
uses a tailored simultaneous width statistic, and the block-incidence
argument translates it into distinct Polynomial roots.  The certificate
machinery relates the displaced-zero comparison and coefficient-order
statistics to the reduced-height literature~\cite{coefficient-mass}.

\section{Routing forces distinct real roots}\label{sec:routing}

Fix an order of the $n$ input variables.  After the first $k$ reads, each
input prefix determines a residual Boolean function of the remaining
$n-k$ variables.  Write $N(k)$ for the number of distinct residuals at that
level and $N'(k)$ for the nonconstant ones.  Call a function $g(y_1,\ldots,y_m)$
\emph{first-essential} if it depends on $y_1$ and is not a function of $y_1$
alone, and write $N^*(k)$ for the number of distinct first-essential residuals
at level $k$.

\begin{lemma}[Maximal-width tables]\label{lem:maximal-width}
For every sufficiently large $n$, there is an $n$-input truth table such
that, under every variable order, the level $k=n-\lceil\log_2n\rceil-1$
satisfies
\begin{equation}\label{eq:residual-width}
  N^*(k+1)=\Omega(T/\log T).
\end{equation}
\end{lemma}

\begin{proof}
Choose a truth table uniformly at random.  Put
$\ell=\lceil\log_2 n\rceil+1$ and expose the residuals after
$k=n-\ell$ variables.  For any fixed variable order, the
$q=2^k=\Theta(T/n)$ residuals occupy disjoint blocks of the truth table and
are therefore independent uniform strings of length
$R=2^\ell\in[2n,4n)$.  Their $2q$ children at level $k+1$ are independent
uniform strings of length $R/2\in[n,2n)$.  Two prescribed children agree
with probability $2^{-R/2}$.

If $N'(k+1)<q/2$, then the $2q$ children occupy fewer than $q/2$ distinct
nonconstant values, so the child occupancy experiment has at least
$3q/2-2\ge q/4$ repetitions (constant children count as repetitions against
either constant string).  Put
$B=2^{R/2}\ge T$ and $t=\lceil q/4\rceil$.  If at least $t$ repetitions
occur, some $t$ children each equal an earlier child.  Choosing those children
and one earlier witness for each gives
\[
  \Pr\{\text{at least }q/4\text{ repetitions}\}
  \le \binom{2q}{t}(2q/B)^t
  \le \left(\frac{4e q^2}{tB}\right)^t
  =\exp(-\Omega(q\log n)),
\]
because $q/B=O(1/n)$.  Constant strings are two ordinary occupancy
bins, so the same bound includes them.

A child fails to be first-essential if its two halves are equal, with
probability $2^{-R/4}$, or if each half is constant, with probability
$4\cdot2^{-R/2}$; so it fails with probability $p\le2^{1-R/4}\le2^{1-n/2}$,
independently across children.  With $t'=\lceil q/8\rceil$,
\[
  \Pr\{\text{at least }q/8\text{ children fail}\}
  \le\binom{2q}{t'}p^{t'}\le(16ep)^{t'}=\exp(-\Omega(qn)).
\]
Outside both events, $N^*(k+1)\ge q/2-q/8=3q/8$.
There are only $n!$ input orders, and
$\log(n!)=O(n\log n)=o(q\log n)$, so with positive probability every order has
$N^*(k+1)\ge 3q/8$.  Since $q=\Theta(T/\log T)$, such a table satisfies
\eqref{eq:residual-width}.
\end{proof}

We call any table satisfying Lemma~\ref{lem:maximal-width}
\emph{maximal-width}.

Call an instruction position \emph{routing} if both of its possible
successors occur on reachable executions, and
let $\mroute$ be the number of routing real-instruction positions.  For
a fixed level $k$, let $D^*_k$ be the number of distinct cursors $c$ such that
some execution performs its $k$th read at $c$ with a first-essential residual.

An execution that halts before its $k$th read has a constant residual.
Otherwise the residual after $k$ reads is determined by the cursor of the
$k$th read together with the register, and a read of an input bit leaves the
register at $48$ or $49$; hence
\begin{equation}\label{eq:residual-cursor}
  N^*(k)\le 2D^*_k .
\end{equation}
Follow an execution back from its $k$th read.  If its path contains a routing
position, let $(r,s)$ be the last one and the successor taken there.  From
$r$ on every real position is non-routing, so its successor is forced, while
reads and register instructions do not branch: every execution that takes
$s$ at $r$ continues along one fixed cursor sequence until the next routing
position or halt.  Let $c_1,\ldots,c_J$ be the reads on this sequence, so
that the $k$th read is some $c_j$.  The index $j$ depends on how many reads
preceded $r$, which may vary between executions, but $c_j$ is among the last
two reads:

\begin{lemma}[Essential reads]\label{lem:essential}
If $k\le n-2$ and the residual after the $k$th read at $c_j$ is
first-essential, then $j\ge J-1$.
\end{lemma}

\begin{proof}
Suppose $j\le J-2$.  After $c_j$ the execution follows the fixed sequence and
reads $y_1=x_{k+1}$ at $c_{j+1}$ and $y_2$ at $c_{j+2}$, with no routing
position in between.  The read at $c_{j+2}$ overwrites the register, so the
state after it does not depend on $y_1$.  Between the two reads $y_1$ enters
only the register, which reaches the cursor only through routing positions,
of which there are none; so it affects only the output emitted there, say
$w(y_1)$.  The run's output is $w(y_1)v(y_2,\ldots)$.  If $w(0)$ and $w(1)$
are empty the residual ignores $y_1$.  Otherwise some $w(y_1)$ is already the
program's single output character, so $v$ is empty for every continuation,
each $w(y_1)$ is one character, and the residual is a function of $y_1$
alone.  A fault between the two reads, such as $0^{-1}$, would occur on a
reachable input, so a program computing the table has none.
\end{proof}

If the path contains no routing position, the same holds for the fixed
sequence from the initial cursor.  Thus every cursor counted by $D^*_k$ is
one of the last two reads before a routing position and successor, or
before the first routing position, so
\begin{equation}\label{eq:basic-routing}
  D^*_k\le 2+4\mroute,
  \qquad\text{and hence}\qquad
  N^*(k)\le 4+8\mroute\qquad(k\le n-2).
\end{equation}

It is not enough to count positions: repeated roots are legal.  We therefore
bound routing positions by the number of distinct root values.

\begin{lemma}[Block-incidence lemma]\label{lem:block}
If a decoded program contains $L$ distinct real instruction-root values,
then
\[
  \mroute\le 3L
  \quad\text{and hence}\quad
  N^*(k)\le 24L+4\qquad(k\le n-2).
\]
\end{lemma}

\begin{proof}
Root sorting places every occurrence of one value in a contiguous block.
Draw each matched control-flow pair as an arc above the instruction line and
contract each block to a vertex.  The arcs are noncrossing, so after parallel
edges are identified the block-incidence graph is a simple outerplanar graph.
There are no loops: an opener and closer have different instruction codes,
whereas equal root values decode to the same code.  It is also bipartite,
with opener blocks on one side and closer blocks on the other.  An outerplanar graph on $L\ge2$ vertices has at most $2L-3$ edges, by
the standard bound applied to each component.

At most one opening instruction routes in each block.  Indeed, after the
first opener in a block, execution can enter the next equal opener only by
falling through its predecessor or by a loop edge landing immediately after
an equal test.  Both entries certify the same truth value of the unchanged
register, so the later opener cannot take both successors.

Charge a routing closer to the incidence between its block and its partner's
opener block.  Noncrossing implies that the closers with a fixed partner block
form one interval inside their own block: nesting would place an
intervening closer's partner strictly between two openers of the partner
block, hence inside that contiguous block.  A closer tests its partner
opener's condition, and a closer whose partner does not loop never jumps.  After the first,
entry is either fallthrough with the common condition false or a skip landing
after the preceding closer, again with that condition false; all intervening
positions are equal real tests and do not change the register.  Thus each
incidence is charged at most once.  For $L\ge2$ there are at most $L$ routing openers
and at most $2L-3$ routing closers; for $L=1$ the displayed weaker bound is
immediate.  Substitute this in \eqref{eq:basic-routing}.
\end{proof}

\begin{corollary}\label{cor:roots}
Every Polynomial program for a maximal-width truth table has
$L=\Omega(T/\log T)$ distinct positive real instruction roots.
\end{corollary}

\begin{proof}
The level $k+1=n-\lceil\log_2n\rceil$ of Lemma~\ref{lem:maximal-width} is
at most $n-2$ once $n\ge3$.  Combine \eqref{eq:residual-width} with
Lemma~\ref{lem:block}.
\end{proof}

\begin{remark}[Why essential residuals]\label{rem:routing-gap}
Counting all read cursors is not enough when the read count varies.  The
decoded list
\[
[5,1],[5],[0,2],[48,2],[2],[0,2],[0,2],[0,2],[0,2],[0,1]
\]
has $\mroute=1$ and four distinct cursors at its fourth read, against
the $1+2\mroute=3$ that one fixed read index per cursor would allow:
its loop reads an input-dependent number of bits, so the routing-free suffix
lands the fourth read at different offsets.  By
Lemma~\ref{lem:essential} only the last two reads of that suffix can carry
first-essential residuals, and here none does: the list prints the last bit
it reads, so every residual ignores its next input or depends on it alone.  Programs that read exactly $n$
bits have at most $1+2\mroute$ distinct $k$th-read cursors, counting
all of them and not only the first-essential ones, but the lower bound
does not need it.
\end{remark}

\section{Coefficient mass from prescribed roots}\label{sec:mass}

The exponential-sum certificate machinery yields the following
degree-independent statement~\cite[Corollary~3.4]{coefficient-mass}.

Write
\[
  \Lambda(F)=\sum_{j:f_j\ne0}\log^+|f_j|.
\]

\begin{theorem}[Coefficient-mass theorem]\label{thm:mass}
Let
\[
  P(x)=\prod_{i=1}^{L}(x-r_i),
  \qquad 2\le r_1\le\cdots\le r_L,
\]
with repetitions allowed, and let \(F(x)=\sum_{j=0}^{D}f_jx^j\in\R[x]\) be a
nonzero multiple of \(P\) with \(|f_D|\ge1\).  Then
\begin{equation}\label{eq:mass-general}
  \Lambda(F)
  \ge (L+1)\log|f_D| + \sum_{u=0}^{L-1}
       \sum_{i=u+1}^{L}\log(r_i-1),
\end{equation}
where \(\log^+x=\max(0,\log x)\).
\end{theorem}

For the first \(L\) primes, or for increasing roots bounded below by them,
the right side is at least \((1/2+o(1))L^2\log L\): each
\(\log(r_i-1)\) occurs \(i\) times, and \(r_i\ge i+1\).  The lower bound is
asymptotically sharp: \(\Lambda(P_L)\le(1/2+o(1))L^2\log L\) for
\(P_L=\prod_{i=1}^L(x-p_i)\).  The theorem applies to every real multiple,
independently of its degree and additional roots.

\begin{corollary}[Polynomial lower bound]\label{cor:lower}
Every Polynomial source for a maximal-width truth table has
\[
  |f|=\Omega(T^2/\log T).
\]
\end{corollary}

\begin{proof}
By Corollary~\ref{cor:roots}, the program contains
\(L=\Omega(T/\log T)\) distinct positive real instruction roots.  They are
distinct integers \(\ge2\), so after sorting \(r_i\ge i+1\), and the weighted
sum in Theorem~\ref{thm:mass} is at least
\(\sum_{i\le L}i\log i=(1/2+o(1))L^2\log L\).  The source polynomial
is an integer multiple of their distinct-root product and its leading
coefficient has absolute value at least one.  Theorem~\ref{thm:mass} gives
\[
 \Lambda(f)=\Omega(L^2\log L)
           =\Omega\!\left(\frac{T^2}{\log T}\right).
\]
Since integer coefficients satisfy
$M(f)\ge\Lambda(f)/\log 10$ and $|f|\ge M(f)$, the result follows.
\end{proof}

\begin{remark}[Repeated roots]
Repeated real roots do not weaken Corollary~\ref{cor:lower}: the
block-incidence lemma already forces the stated number of distinct values,
and Theorem~\ref{thm:mass} is stated for roots counted with multiplicity.
\end{remark}

\section{The matching construction}\label{sec:upper}

The residual DAG of a truth table has at most
\[
  s_k=\min\bigl(2^k,2^{2^{n-k}}\bigr)
\]
states after $k$ reads: the first term counts input prefixes, and the second
counts Boolean residuals of the remaining width.  Put
$h=\lfloor\log_2 n\rfloor-1$.  The levels with $n-k\le h$ contribute
$O(\sqrt T)=O(T/n)$ states; the earlier levels form a geometric tail
of size $O(T/2^h)=O(T/n)$.  Hence
\[
  S\coloneqq\sum_k s_k=O(T/\log T).
\]

Label the states at every level by $0,\ldots,s_k-1$.  Enter level $k$ with
register $-j$ when state $j$ is active.  Its blocks are scanned in label
order: block $j$ increments the register and opens an ``if positive'', so
exactly block $j$ fires.  That block reads the next bit, overwriting the
register with $48+b$.  If its child labels are $z,o$, multiply by $z-o$ and
add $-(z+r)-48(z-o)$, where $r$ is the number of blocks remaining at that
level.  This sends $48$ to $-(z+r)$ and $49$ to $-(o+r)$.  The skipped blocks
each increment once, leaving $-z$ or $-o$ at the next level.  Equal children
instead floor-divide $48+b$ by $50$, obtaining zero, then add $-(z+r)$.  At
the leaf level the unique active block adjusts the register to $48$ or $49$,
prints, and leaves the register negative.

Thus induction on the levels proves that
the machine follows exactly the residual selected by the input and prints its
table value.  Each state block uses at most six instructions, its brackets
are local and therefore properly matched, and every operand is $O(S)$.  An
update whose operand would be zero is omitted: the language spells a zero
operand as I/O, and dropping it leaves the register unchanged, which is what
a zero update does.

Encode the resulting $m=O(S)$ instructions in order, using one fresh prime
per instruction.  Real instructions give linear factors and register
instructions conjugate quadratic factors; multiplying them gives an integer
polynomial decoded back to that list.  Instruction exponents are bounded
constants.  The prime number theorem~\cite{rs} gives the standard estimate
$p_m=O(m\log m)$ for the largest assigned prime.  Every
linear or quadratic factor therefore has
coefficient $\ell^1$ norm $m^{O(1)}$.  Submultiplicativity gives coefficient
norm $m^{O(m)}$ for their product, so each expanded coefficient has
$O(m\log m)$ decimal digits.  There are at most $2m+1$ coefficients, while
exponents and separators contribute only $O(m\log m)$.  Therefore
\[
  |f|=O(m^2\log m)
      =O\!\left(\frac{T^2}{\log T}\right).
\]
Together with Corollary~\ref{cor:lower}, this proves
Theorem~\ref{thm:main}.

\section{Explicit constants}\label{sec:constants}

Theorem~\ref{thm:main} is a statement of order.  This section tracks the
constants through both halves.  Write
\[
  C_P(n)=\max_{t\in\{0,1\}^{T}}\ \min\{|f|:f\text{ computes }t\},
  \qquad T=2^n,
\]
and normalize by $T^2/n=T^2/\log_2T$.  Two small sharpenings of the routing
argument are proved first; neither changes the order.

\begin{lemma}[One residual below the last read]\label{lem:essential-sharp}
In the setting of Lemma~\ref{lem:essential}, let $k\le n-2$ and suppose the
$k$th read is at $c_{J-1}$.  If the residual after it is first-essential,
then that residual does not depend on the bit just read.  Consequently
\[
  N^*(k)\le 3+6\mroute\qquad(k\le n-2).
\]
\end{lemma}

\begin{proof}
Let $b$ be the bit read at $c_{J-1}$.  From there the execution follows the
fixed routing-free sequence to $c_J$ and performs no read in between.  At
$c_J$ it reads $y_1$, and this read overwrites the register.  So the state
after $c_J$ is $(c_J,48+y_1)$ whatever $b$ was, and the output after the
$k$th read has the form $w(b)\,v(y_1,y_2,\ldots)$.  Here $w(b)$ is what is
printed between the two reads, and $v$ does not depend on $b$.  A fault
between the reads would lie on a reachable input.  If $w(b)$ is nonempty it
is the program's single output character, and the residual is constant.  A
first-essential residual therefore has $w(b)$ empty and equals $v$, for
either value of $b$.

Now count cursors by segment.  There are at most $2\mroute+1$ segments:
one for each routing position and successor taken there, and one from the
initial cursor.  By Lemma~\ref{lem:essential}, each segment contributes
only the two cursors $c_J$ and $c_{J-1}$.  The last read $c_J$ carries at
most two first-essential residuals, one for each register value $48$ or
$49$, as in~\eqref{eq:residual-cursor}.  By the first paragraph, $c_{J-1}$
carries at most one.  Hence $N^*(k)\le3(2\mroute+1)$.
\end{proof}

\begin{lemma}[Bipartite block incidence]\label{lem:bipartite}
If a decoded program has $L\ge2$ distinct real root values, then
$\mroute\le 2L-2$, and hence $N^*(k)\le 12L-9$ for $k\le n-2$.
\end{lemma}

\begin{proof}
Let $a$ and $b$ be the numbers of opener blocks and closer blocks, so
$L=a+b$.  Let $G$ be the block-incidence graph of Lemma~\ref{lem:block},
with $E$ edges.  That proof bounds the routing positions by $a+E$.  $G$ is
simple, outerplanar and bipartite, with parts of sizes $a$ and $b$.  We
claim that $E\le a+2b-2$ whenever $E\ge1$.

First take a $2$-connected block of $G$ with $V_i\ge3$ vertices.  Its outer
face is a Hamiltonian cycle, and this cycle alternates between the two
parts, so $a_i=b_i$.  Every inner face has length at least~$4$.  Euler's
formula then gives $2E_i\ge4(F_i-1)+V_i$ and hence
$E_i\le(3V_i-4)/2=a_i+2b_i-2$.  A bridge has $a_i=b_i=1$ and $E_i=1$, so it
satisfies the same bound.

In a connected component, cut vertices are counted once in each block that
contains them.  The total overcount is $x+y$, where $x$ comes from openers
and $y$ from closers, and it equals the number of blocks minus one.
Summing over the blocks gives
\[
  E\le(a+x)+2(b+y)-2(x+y+1)=a+2b-2-x\le a+2b-2 .
\]
Components without edges only increase the right side, and further
components each subtract another~$2$.  If $E=0$, then $\mroute\le a\le L$.
Otherwise $\mroute\le a+E\le2(a+b)-2$.  Substitute this in
Lemma~\ref{lem:essential-sharp}.
\end{proof}

The count $a+E=2V-2$ is attained, for instance by a chain of $4$-cycles.
So the graph step is sharp.  It is the passage from segments to cursors
that is lossy (Remark~\ref{rem:bracket}).

\begin{proposition}[Explicit bracket]\label{prop:bracket}
For $2\le j\le n$ put $q_j=2^{n-j}$ and
$E_j=2^{2^j}-2^{2^{j-1}}-2$, and let
\[
  N_n=\max_{2\le j\le n}\min(q_j,E_j),\qquad
  s_k=\min\bigl(2^k,2^{2^{n-k}}\bigr),\qquad
  S_n=\sum_{k<n}s_k .
\]
Let $k_1$ be the largest $k$ with $s_k=2^k$, let $R_n=\sum_{k_1<k<n}s_k$,
and let $\Phi_n=\tfrac12\bigl(311S_n^2-(10S_n-2R_n)^2\bigr)$.  Then:
\begin{enumerate}
\item for every $n\ge6$ (where $N_n\ge10$), with $L_n=\lceil(N_n+9)/12\rceil$,
\[
  C_P(n)\ \ge\ \frac1{\log10}\sum_{i=1}^{L_n}i\log i
        \ \ge\ \frac{L_n^2(2\log L_n-1)}{4\log10};
\]
\item every $n$-input table has a program with
\[
  |f|\le\bigl(1+O(1/n)\bigr)\,n\log_{10}2\cdot\Phi_n .
\]
\end{enumerate}
Put $j_1=\min\{j:2^j+j\ge n\}$, $\nu_n=n/2^{j_1}$, write
$n=2^{j_1-1}+j_1-1+s$ with $s\ge1$, and let $u=2^{-s}$.  Then
\begin{equation}\label{eq:bracket-n}
  \Bigl(1-O\bigl(\tfrac{\log n}n\bigr)\Bigr)\frac{\log_{10}2}{288}\,\nu_n^2
  \ \le\ \frac{C_P(n)\,n}{T^2}\ \le\
  \bigl(1+o(1)\bigr)\log_{10}2\;\nu_n^2\bigl(422+924u+494u^2\bigr),
\end{equation}
and consequently
\[
  \frac{\log_{10}2}{1152}\le\liminf_{n\to\infty}\frac{C_P(n)\,n}{T^2}
  \le\frac{211}{2}\log_{10}2,
  \qquad
  \frac{\log_{10}2}{288}\le\limsup_{n\to\infty}\frac{C_P(n)\,n}{T^2}
  \le422\log_{10}2 .
\]
Numerically, $2.61\cdot10^{-4}\le\liminf\le31.8$ and
$1.05\cdot10^{-3}\le\limsup\le127.1$.
\end{proposition}

\begin{proof}
\emph{Lower bound.}  Reads consume the input stream in order, so the
residuals that matter are those of the natural variable order.  The table
can therefore be chosen outright, and Lemma~\ref{lem:maximal-width}'s union
bound over orders is not needed.  Fix the maximizing $j$ and $K=n-j\le n-2$.
The level-$K$ residuals are the $q_j$ consecutive blocks of length $2^j$
and may be prescribed independently.  Exactly $E_j$ strings of that length
are first-essential: we exclude the $2^{2^{j-1}}$ strings with equal
halves, and the two further strings whose halves are distinct constants.
Prescribing $\min(q_j,E_j)$ distinct first-essential blocks gives a table
with $N^*(K)=N_n$.

A program with $L\le1$ distinct real values has $\mroute\le1$ and so
$N^*(K)\le9<N_n$; hence $L\ge2$, and Lemma~\ref{lem:bipartite} gives
$L\ge L_n$ distinct real instruction roots.
They are distinct integers at least~$2$, so $r_i-1\ge i$ after sorting.
Theorem~\ref{thm:mass} and $|f|\ge M(f)\ge\Lambda(f)/\log10$ give the first
inequality.  The second is
$\sum_{i\le L}i\log i\ge\int_1^Lx\log x\,dx$.

For the asymptotic form, note that $2^{j}\ge n-j$ exactly when $j\ge j_1$.
So $q_{j_1}\le 2^{2^{j_1}}$, and $\min(q_{j_1},E_{j_1})=q_{j_1}(1+o(1))$.
Also $2^{2^{j_1-1}}=2^{1-s}q_{j_1}\le q_{j_1}$, and $q_j$ halves as $j$
grows.  Hence $N_n=q_{j_1}(1+o(1))=\nu_nT/n\,(1+o(1))$.  Then
$\log L_n=n\log2-O(\log n)$, which gives the left side
of~\eqref{eq:bracket-n}.

\emph{Upper bound.}  We use the construction of Section~\ref{sec:upper},
with primes shared along runs of nondecreasing decode key, as the generator
does; each prime then carries at most three instructions.  A chained state
block is
\[
  \mathord{+}\!=1,\ \text{if}>0,\ \text{input},\
  \mathord{*}\!=\text{span}\ \text{or}\ /\!/\!=50,\
  \mathord{+}\!=\delta,\ \text{endif}.
\]
Its ten roots have moduli $|a+p^bi|$ or $p^v$, with exponent $b$ (or $v$)
equal to $1$ five times, $2$ three times, and $3$ or $4$ twice; the value
$4$ arises only for a state with equal children.  Omitted instructions
only remove roots, and the leaves add $O(1)$ roots.

For a multiset of positive exponents let $\Phi=\sum_ii\,b_{(i)}$, taken in
increasing order.  Suppose there are $S'$ chained blocks, $Y$ of them with
equal children.  A direct summation gives
$\Phi\le\frac12\bigl(311S'^2-(10S'-2Y)^2\bigr)+O(S')$.

Next we bound $Y$.  Let $w_k$ be the widths and $f_k=s_k-w_k$ the
deficits, and let $Y_k$ count the equal-children states at level $k$.  For
$k<k_1$ we have $s_{k+1}=2s_k$ and $w_{k+1}\le2w_k-Y_k$, so
$f_{k+1}\ge2f_k+Y_k$, and $f_0=0$.  Hence
$S'\le S_n-\sum_{k<k_1}Y_k$.  For $k\ge k_1$ we have
$Y_k\le w_{k+1}$, so $\sum_{k\ge k_1}Y_k\le R_n+2$.  Put
$g=311S'^2-(10S'-2Y)^2$.  It increases in $Y$ for $Y\le S'$.  Along
$S'=S_n-t$, $Y=R_n+t$ its derivative in~$t$ is
$2(-191S_n+167t-24R_n)<0$.  Hence $\Phi\le\Phi_n+O(S_n)$.

Every nonleading coefficient is an elementary symmetric function of the
$D\le10S_n+O(1)$ roots, and all moduli exceed~$1$, so
\[
  \Lambda(f)\le\sum_{k}\log\binom Dk+\sum_{i}i\log|\rho|_{(i)} ,
\]
the second sum taken over the moduli in increasing order.

Write $\log|\rho|\le b\log P+\varepsilon_\rho$, where $P$ is the largest
prime and $\varepsilon_\rho=\log(1+|a|/p^b)$.  By the rearrangement
inequality the second sum is at most $\Phi\log P+D\sum_\rho\varepsilon_\rho$.

The operands satisfy $|a|\le A=50S_n+n+60$.  A root with $p\ge A$ has
$\varepsilon_\rho\le\log2$.  At most $6\pi(A)=O(S_n/\log S_n)$ roots have
$p<A$, and each has $\varepsilon_\rho\le\log(1+A)$.  So
$D\sum\varepsilon_\rho=O(S_n^2)$.  Also $\sum_k\log\binom Dk=O(S_n^2)$.
There are $O(S_n)$ primes, so~\cite{rs} gives $\log P=n\log2+O(1)$, while
$\Phi_n\ge50S_n^2$.  Therefore $\Lambda(f)\le(1+O(1/n))\,n\log2\cdot\Phi_n$.

Rendering adds at most $D+1$ digits and $O(D\log D)$ further characters.
Dividing by $\log10$ proves the second item.

Finally, $k_1=n-j_1$, so
$S_n=2q_{j_1}-1+R_n$.  The sum $R_n$ is dominated by its first term, so
$R_n=2^{2^{j_1-1}}(1+o(1))=2u\,q_{j_1}(1+o(1))$.  With $q_{j_1}=\nu_nT/n$,
substitution gives
$\Phi_nn^2/T^2=\nu_n^2(422+924u+494u^2)(1+o(1))$.

Along $n=2^{j}+j$ we have $\nu_n\to1$ and $u\to0$.  Along
$n=2^{j-1}+2j$ we have $\nu_n\to\frac12$ and $u\to0$.  For every $n$,
$\nu_n>\frac12$, and $\nu_n\le1+j_1/2^{j_1}$.  The bound
$\nu_n^2(422+924u+494u^2)$ tends to $422$ along the first sequence and to
$105.5$ along the second.  It never exceeds $422+o(1)$: large $\nu_n$ needs
$s\approx 2^{j_1-1}$, and at $s=1$ its value is about
$\frac14\cdot1007.5$.  This proves the limit statements.
\end{proof}

\begin{remark}[What the bracket says]\label{rem:bracket}
At every $n$ the two sides of~\eqref{eq:bracket-n} differ by the factor
$288(422+924u+494u^2)$.  This is $121{,}536$ for most $n$ ($u\to0$) and at
most $290{,}160$ (at $s=1$).  So Proposition~\ref{prop:bracket} fixes the
order, but it does not fix the constant.  At $u\to0$ the factor splits as
\[
  121{,}536\ =\ 12^2\ \times\ 4\ \times\ 2\cdot105.5 .
\]
\begin{itemize}
\item \emph{$12^2$: routing.}  The lower bound credits one distinct real
root to every $12$ first-essential residuals at one level.  The
construction spends two per state.  Both per-cursor caps in
Lemma~\ref{lem:essential-sharp} are attained, and so is the graph bound in
Lemma~\ref{lem:bipartite}.  The loss comes from combining them.
\item \emph{$4=(S_n/q_{j_1})^2$: levels.}  The lower bound uses a single
level.  The construction pays for all of them.  Counting cannot add levels,
because one cursor may serve residuals at several levels when the read
count varies.
\item \emph{$2\cdot105.5$: roots the certificate does not see.}  The mass
constant $\frac12$ of Theorem~\ref{thm:mass} is sharp: the product of the
first $L$ primes attains it~\cite[Corollary~3.6]{coefficient-mass}.  The
logarithms satisfy $\log L\sim\log P\sim\log T$.  But the certificate
charges only real roots.  In the construction, $80\%$ of the degree
consists of the complex register roots, and it pays
$\Phi_n/S_n^2\le105.5$ against $\frac12$.  Random tables measure about
$76$.
\end{itemize}
In particular, the order-statistic bound for arbitrary free positions
(\cite[Lemma~3.1]{coefficient-mass}) is already the sharp form of the mass
step.  No sharpening of that step can move the bracket by more than the
factor it already has, which is exact.  What remains is combinatorial: a
routing lemma with a constant near the construction's, and a coefficient
mass bound for the complex instruction roots.
\end{remark}

\section{Scope and reproducibility}

The result concerns source length, not the running time of polynomial
factorization or of source generation.  It allows arbitrary degrees,
arbitrary operand signs, repeated roots, and arbitrary unrecognized cofactor
roots.  It is an existence theorem for the uncapped residual-DAG
construction.  Implementation resource limits do not affect the existence
result.

Executable checks cover the finite block-counting identities and repeated-root
cases used here.  They support the proof, but no bounded computation is used
as a premise.

\paragraph{Data availability.}
No datasets were generated or analyzed in this study.

\begin{thebibliography}{9}
\small
\bibitem{polynomial-spec}
Maedhros777,
\newblock Polynomial,
\newblock Esolang Wiki,
\newblock \href{https://esolangs.org/w/index.php?title=Polynomial\&oldid=192583}
{permanent revision 192583}.

\bibitem{rs}
J.~B. Rosser and L.~Schoenfeld,
\newblock Approximate formulas for some functions of prime numbers,
\newblock \emph{Illinois Journal of Mathematics} 6 (1962), 64--94,
\newblock \href{https://doi.org/10.1215/ijm/1255631807}
{doi:10.1215/ijm/1255631807}.

\bibitem{bryant}
R.~E. Bryant,
\newblock Graph-based algorithms for Boolean function manipulation,
\newblock \emph{IEEE Transactions on Computers} 35 (1986), 677--691,
\newblock \href{https://doi.org/10.1109/TC.1986.1676819}
{doi:10.1109/TC.1986.1676819}.

\bibitem{wegener}
I.~Wegener,
\newblock The size of reduced OBDD's and optimal read-once branching programs
for almost all Boolean functions,
\newblock \emph{IEEE Transactions on Computers} 43 (1994), 1262--1269,
\newblock \href{https://doi.org/10.1109/12.324559}
{doi:10.1109/12.324559}.

\bibitem{gps}
C.~Gr\"opl, H.~J. Pr\"omel, and A.~Srivastav,
\newblock Ordered binary decision diagrams and the Shannon effect,
\newblock \emph{Discrete Applied Mathematics} 142 (2004), 67--85,
\newblock \href{https://doi.org/10.1016/j.dam.2003.02.003}
{doi:10.1016/j.dam.2003.02.003}.

\bibitem{coefficient-mass}
B.~Pham,
\newblock Displaced-zero certificates and the coefficient mass of polynomial
multiples,
\newblock manuscript in preparation; Corollary~3.4 (Logarithmic mass).

\end{thebibliography}

\end{document}
